\chapter{Introduction}

GitHub Actions is a CI/CD system driven by events in your repository.
A workflow is a YAML file stored in \code{.github/workflows}.
It reacts to events, runs jobs on runners, and produces outputs you can trust.

This book assumes you can read YAML and you know what a CI pipeline should do.
The goal is to build automation that is predictable, secure, and easy to change.

\section{Who this book is for}
This book is for engineers who own real repositories:
\begin{itemize}[nosep]
  \item teams with multiple services or packages
  \item projects that need repeatable builds and tests
  \item repos with a mix of contributors and release cadence
\end{itemize}

If you are entirely new to CI, read this alongside a short CI primer.
If you already use GitHub Actions, use this as a guide to harden and scale workflows.

\section{What you will build}
You will build workflows that cover:
\begin{itemize}[nosep]
  \item lint, type-check, and tests on every pull request
  \item caching to keep builds fast and deterministic
  \item artifacts for reports and binaries
  \item composite actions for reuse inside a repo
  \item reusable workflows for reuse across repos
  \item secure permissions and secret handling
  \item release and deployment pipelines that are auditable
\end{itemize}

\section{How to read this book}
Each chapter introduces a narrow concept and a small, runnable example.
Copy the examples into a sandbox repository and adjust them to your stack.
Avoid premature complexity; only add features you can explain to a teammate.

\begin{patternbox}{Pattern: treat workflows as product code}
If a workflow is hard to read, it will be hard to maintain.
Keep naming consistent, add comments only where needed, and prefer scripts over inline logic.
\end{patternbox}

\section{Repository layout}
Most examples assume this structure:
\begin{itemize}[nosep]
  \item \code{.github/workflows/} for workflow files
  \item \code{scripts/} for repeatable project commands
  \item \code{docs/} for generated reports or artifacts
\end{itemize}

\section{A first workflow}
This example runs on pushes to \code{main} and on pull requests.
It checks out code and runs a script.

\begin{lstlisting}[language=YAML]
name: CI (basic)

on:
  push:
    branches: ["main"]
  pull_request:

permissions:
  contents: read

jobs:
  hello:
    runs-on: ubuntu-latest
    steps:
      - uses: actions/checkout@v4
      - name: Say hello
        run: |
          echo "Hello from GitHub Actions"
\end{lstlisting}

\section{A realistic first pipeline}
Most teams start with lint and tests, then add cache and artifacts.

\begin{lstlisting}[language=YAML]
name: CI (lint + test)

on:
  pull_request:
  push:
    branches: ["main"]

permissions:
  contents: read

jobs:
  test:
    runs-on: ubuntu-latest
    steps:
      - uses: actions/checkout@v4

      - name: Install
        run: ./scripts/install

      - name: Lint
        run: ./scripts/lint

      - name: Test
        run: ./scripts/test
\end{lstlisting}

\section{Common early mistakes}
\begin{itemize}[nosep]
  \item Keeping the workflow logic in YAML instead of scripts.
  \item Running privileged steps in pull request workflows.
  \item Ignoring caching until the pipeline is already slow.
\end{itemize}

\section{Scope and non-goals}
This book focuses on practical workflow design, not on every GitHub Actions feature.
We do not cover marketplace action authoring in depth or advanced runner orchestration.
Use the patterns here as a foundation, then layer in advanced features as you need them.

\section{A minimal scripts layout}
Move logic out of YAML into small scripts so it can be run locally.
Keep the scripts boring and composable.

\begin{lstlisting}[language=YAML]
scripts/
  install
  lint
  test
  build
\end{lstlisting}

\begin{checklistbox}{Checklist: keep the first workflow boring}
\begin{itemize}[nosep]
  \item Add \code{permissions} explicitly.
  \item Prefer small steps with clear names.
  \item Keep \code{run} scripts short; move logic to repo scripts as soon as it grows.
  \item Treat the workflow as a thin orchestrator, not the place for business logic.
\end{itemize}
\end{checklistbox}
